\documentclass[UTF8,a4paper,12pt]{ctexart}
\usepackage{geometry}
\usepackage{xcolor}
\usepackage{hyperref}
\usepackage{listings}
\usepackage{longtable}
\usepackage{array}
\usepackage{enumitem}

\geometry{left=2.6cm,right=2.6cm,top=2.5cm,bottom=2.5cm}
\hypersetup{
  colorlinks=true,
  linkcolor=blue,
  urlcolor=blue
}
\lstset{
  basicstyle=\ttfamily\small,
  breaklines=true,
  frame=single,
  columns=fullflexible,
  backgroundcolor=\color{gray!8}
}
\setlist[itemize]{itemsep=4pt}
\setlist[enumerate]{itemsep=4pt}

\title{哈夫曼树创建过程交互演示说明文档}
\author{数据结构课程作业}
\date{2026年5月}

\begin{document}

\maketitle

\section{项目概述}

本项目围绕“哈夫曼树创建”这一数据结构主题，制作了一个可部署到 GitHub Pages 的交互式网页。用户可以输入字符与权值，点击按钮逐步观察哈夫曼树的构建过程：初始化优先队列、取出两个最小节点、合并新节点、重新入队，直到队列中只剩根节点。构建完成后，网页会根据“左分支为 0，右分支为 1”的规则生成哈夫曼编码表。

项目最终包含三个主要成果：

\begin{itemize}
  \item 交互式网页：\texttt{index.html}、\texttt{styles.css}、\texttt{app.js}。
  \item 说明文档源文件：\texttt{report.tex}。
  \item 说明文档 PDF：\texttt{report.pdf}。
\end{itemize}

部署后的访问地址规划为：

\begin{lstlisting}
https://ccycyc6.github.io/ds_work/
\end{lstlisting}

\section{哈夫曼树创建思路}

哈夫曼树是一种带权路径长度最短的二叉树，常用于压缩编码。构建过程可以用优先队列描述：

\begin{enumerate}
  \item 将每个字符及其权值看作一个叶子节点。
  \item 将所有叶子节点按权值从小到大放入优先队列。
  \item 每次从队列中取出权值最小的两个节点。
  \item 创建一个新的父节点，其权值等于两个子节点权值之和。
  \item 将两个节点分别作为新节点的左孩子和右孩子。
  \item 把新节点重新放入优先队列。
  \item 重复上述过程，直到队列中只剩一个节点，该节点就是哈夫曼树的根。
\end{enumerate}

网页演示中采用稳定排序规则：如果两个节点权值相同，则按照它们进入队列的先后顺序排序。这样可以保证同一组输入每次得到的演示结果一致。

\section{网页功能说明}

\subsection{输入与初始化}

用户输入格式为：

\begin{lstlisting}
A:5, B:9, C:12, D:13, E:16, F:45
\end{lstlisting}

每一项由字符、冒号和正整数权值组成，多项之间可以使用逗号或换行分隔。点击“初始化”后，程序会检查输入是否为空、字符是否重复、权值是否合法，然后生成初始优先队列。

\subsection{逐步构建}

点击“下一步”时，页面会依次展示以下动作：

\begin{itemize}
  \item 高亮当前取出的两个最小节点。
  \item 显示两个节点合并后的新父节点。
  \item 更新优先队列。
  \item 在 SVG 区域绘制当前树结构。
  \item 在步骤说明中记录本次操作。
\end{itemize}

\subsection{编码生成}

当队列中只剩一个根节点时，构建结束。程序从根节点开始遍历整棵树，左边路径记为 0，右边路径记为 1，最终得到每个叶子字符的哈夫曼编码。

\section{关键实现}

\subsection{算法设计}

核心逻辑位于 \texttt{app.js}。程序先将输入解析为叶子节点数组，然后通过 \texttt{buildSteps} 函数预先生成所有演示步骤。每个步骤都保存当前队列、当前树根、被选中的节点以及步骤说明。这样页面渲染时只需要根据当前步骤更新视图，交互逻辑比较清晰。

核心伪代码如下：

\begin{lstlisting}
queue = sort(leaves)
while queue.length > 1:
    first = queue.popMin()
    second = queue.popMin()
    parent = new Node(first.weight + second.weight)
    parent.left = first
    parent.right = second
    queue.push(parent)
    queue = sort(queue)
root = queue[0]
\end{lstlisting}

\subsection{可视化设计}

树形结构使用 SVG 绘制。程序递归遍历当前根节点，为每个叶子节点分配横向位置，再让父节点位于左右子树的中间。边上标注 0 或 1，节点中显示字符或内部节点编号以及权值。取出的节点、合并节点和叶子节点使用不同颜色区分。

\subsection{页面结构}

页面分为五个区域：

\begin{itemize}
  \item 输入与控制区。
  \item 优先队列区。
  \item 树形动画区。
  \item 步骤说明区。
  \item 哈夫曼编码表区。
\end{itemize}

网页采用原生 HTML、CSS 和 JavaScript 实现，不依赖框架和构建工具，因此可以直接部署到 GitHub Pages。

\section{AI 工具使用说明}

本项目使用 AI 辅助完成需求拆解、交互设计、代码生成、文档撰写和部署说明。AI 的作用不是替代算法理解，而是帮助把课程要求转化为可通过公网访问的网页和规范化说明文档。

\subsection{工具介绍}

\begin{longtable}{|p{3cm}|p{10cm}|}
\hline
工具 & 作用 \\
\hline
ChatGPT / Codex & 分析作业要求，制定实现计划，生成网页代码、JavaScript 算法逻辑、CSS 样式和 LaTeX 文档。 \\
\hline
终端命令 & 检查目录、运行算法校验、编译 LaTeX、管理 Git 提交与推送。 \\
\hline
XeLaTeX & 将中文 LaTeX 源文件编译为 PDF。 \\
\hline
GitHub Pages & 托管静态网页，提供可访问的 Web 链接。 \\
\hline
\end{longtable}

\subsection{提示词记录}

以下记录按实际完成项目的阶段整理，保留了关键提示词和 AI 输出目标。

\subsubsection*{第 1 轮：提出作业需求}

用户提示词：

\begin{lstlisting}
设计并制作一个交互式网页或演示动画，直观地展示某种树形数据结构的构建或操作过程，我要做哈夫曼树创建，做成一个可访问的 Web 链接，然后说明文档(PDF)（用latex）：
关键步骤说明：整体思路阐述，解释如何利用 AI 工具分步完成上述动画。
工具介绍、提示词记录：完整、详细记录你与 AI 的交互过程。
\end{lstlisting}

AI 输出目标：确认作业需要同时包含网页动画、Web 链接、PDF 文档和 AI 交互记录。

\subsubsection*{第 2 轮：确认交付方式}

AI 提问要点：

\begin{lstlisting}
你希望最终的“可访问 Web 链接”是哪一种？
哈夫曼树创建动画希望做到哪种交互深度？
PDF 说明文档使用什么语言和记录粒度？
\end{lstlisting}

用户选择：

\begin{lstlisting}
GitHub Pages
逐步构建
中文详细
\end{lstlisting}

AI 输出目标：将项目定位为可部署到 GitHub Pages 的静态站点，并确定交互方式为逐步构建。

\subsubsection*{第 3 轮：形成实施计划}

AI 生成计划摘要：

\begin{lstlisting}
创建 index.html、styles.css、app.js、report.tex、report.pdf 和 README.md。
网页支持输入权值、初始化、下一步、重置、示例数据。
算法按优先队列构建哈夫曼树，完成后生成编码表。
GitHub Pages 访问地址为 https://ccycyc6.github.io/ds_work/。
\end{lstlisting}

AI 输出目标：给出可执行的文件结构、功能范围和部署方式。

\subsubsection*{第 4 轮：开始实施}

用户提示词：

\begin{lstlisting}
PLEASE IMPLEMENT THIS PLAN
\end{lstlisting}

AI 执行目标：

\begin{itemize}
  \item 编写静态网页文件。
  \item 实现哈夫曼树算法和逐步动画。
  \item 编写中文 LaTeX 说明文档。
  \item 编译 PDF。
  \item 初始化 Git 仓库并连接到 GitHub 仓库。
\end{itemize}

\subsection{AI 分步完成动画的方法}

AI 完成动画时采用了分层拆解的方式：

\begin{enumerate}
  \item 先把哈夫曼树创建算法抽象为“状态序列”，每一个状态代表一个演示步骤。
  \item 再设计页面区域，让输入、队列、树形图、说明文字和编码表分别承担不同的信息展示任务。
  \item 然后编写纯 JavaScript 算法函数，保证不依赖页面也能校验构建结果。
  \item 接着使用 SVG 根据当前树节点递归计算坐标并绘制节点和连线。
  \item 最后补充输入校验、按钮状态、错误提示和响应式 CSS。
\end{enumerate}

\section{公网访问与部署}

\subsection{公网访问地址}

本项目通过 GitHub Pages 提供公网访问。仓库启用 Pages 后，可直接打开：

\begin{lstlisting}
https://ccycyc6.github.io/ds_work/
\end{lstlisting}

如果首次打开时显示 404，通常表示 GitHub Pages 尚未启用或部署还未完成，需要在仓库设置中完成 Pages 配置并等待生成。

\subsection{说明文档}

网页中提供了 PDF 文档入口，公网访问页面后可点击“查看说明文档”打开 \texttt{report.pdf}。文档也可通过仓库中的 \texttt{report.pdf} 文件查看。

\subsection{编译 PDF}

使用 XeLaTeX 编译：

\begin{lstlisting}
xelatex report.tex
xelatex report.tex
\end{lstlisting}

\subsection{部署到 GitHub Pages}

将项目推送到仓库：

\begin{lstlisting}
git init
git branch -M main
git remote add origin https://github.com/ccycyc6/ds_work
git add .
git commit -m "Create Huffman tree interactive demo"
git push -u origin main
\end{lstlisting}

然后在 GitHub 仓库中进入 \texttt{Settings -> Pages}，选择 \texttt{Deploy from a branch}，分支选择 \texttt{main}，目录选择 \texttt{/root}。保存后访问：

\begin{lstlisting}
https://ccycyc6.github.io/ds_work/
\end{lstlisting}

\section{功能验收}

面向最终公网页面，项目应满足以下验收标准：

\begin{lstlisting}
https://ccycyc6.github.io/ds_work/
\end{lstlisting}

\begin{itemize}
  \item 可以输入字符权值表并初始化优先队列。
  \item 可以通过“下一步”逐步观察取出最小节点、合并节点和重新入队的过程。
  \item 树形区域能够同步显示节点、边和 0/1 路径标记。
  \item 构建完成后能够生成哈夫曼编码表。
  \item 页面中的“查看说明文档”可以打开 PDF 文档。
\end{itemize}

\section{总结}

本项目将哈夫曼树的抽象构建过程转化为可观察、可交互的公网网页动画。通过逐步按钮，用户可以清楚看到优先队列如何变化、节点如何合并、树如何增长，以及最终编码如何产生。AI 工具主要用于辅助需求拆解、代码实现、文档组织和部署说明，最终形成了可部署、可演示、可提交的课程作业成果。

\end{document}
